\section{An exact polynomial policy for the kernel LP alternative}
\label{sec:exact-polynomial-lp}

The LP interface used by the Euler searches can be made polynomial in the
binary encoding of its rational input while preserving its two exact output
forms.  The maintained polynomial policy keeps a bounded prefix of the existing
Bland simplex implementation and then uses a fully rational, outward-rounded
ellipsoid algorithm.  An explicit determinant gap makes feasibility decisions
exact; a sequence of tight-row decisions recovers an original feasible point;
and support descent recovers the extreme ray required by conflict learning.
This removes the unbounded-simplex-pivot obligation for this interface.  It
does not bound the number of admissibility branches or prove quasi-polynomial
unknot recognition.

The classical basis for polynomial-bit ellipsoid feasibility and exact
primal/dual optimization is due to Khachiyan and its subsequent refinements.
Our primary references are Gr\"otschel, Lov\'asz and Schrijver,
\emph{Geometric Algorithms and Combinatorial Optimization},
Sections~3.1--3.2 (especially Theorem~3.2.1 and Lemmas~3.2.8--3.2.10) and
Sections~6.4--6.5, including Theorem~6.5.17; and G\'acs and Lov\'asz,
``Khachiyan's algorithm for linear programming'',
\emph{Mathematical Programming Studies}~14 (1981), 61--68.\footnote{\raggedright
Author-hosted book:
\url{https://www.zib.de/userpage/groetschel/pubnew/paper/groetschellovaszschrijver1988.pdf}.
Original G\'acs--Lov\'asz publication:
\url{https://doi.org/10.1007/BFb0120921}.}
These references justify the general exact finite-bit framework.  The
particular gap, grid, rounding constants, recovery procedure, and accounting
below are specified and proved here; a real-arithmetic ellipsoid iteration
bound alone would not justify this implementation.

\subsection{Contract, exact input, and bounded alternatives}

For $C\in\mathbb Q^{m\times n}$ and $c\in\mathbb Q^n$, the required alternative is
\begin{align*}
 &Cx=0,\quad x\ge0,\quad c^{\mathsf T}x>0,
 &&\text{or}\\
 &C^{\mathsf T}y\ge c.
\end{align*}
The positive output is a primitive kernel extreme ray, together with its
normalization $\sum_jx_j=1$ and a compatible slack-form basis.  The negative
output gives a multiplier for every original input row, including rows
omitted during rank preprocessing.  The final inequalities are evaluated
exactly on the original rational input.  A finite work allowance may instead
produce \texttt{INCONCLUSIVE}; no incomplete recovery is a certificate.

Choose $r$ independent \emph{original} rows.  Their rational row span is the
whole row space of $C$.  Multiply selected row $i$ by a positive integer
$s_i$, and the objective by a positive integer $s_c$, clearing denominators:
\[
 C_I=\operatorname{diag}(s_i)C_{\rm selected},
 \qquad c_I=s_c c.
\]
Least common multiples of input denominators have polynomial encoding
length.  Let $H\ge1$ bound the absolute integer entries in $C_I$ and $c_I$.
For $n\ge1$, put
\[
 k=\min(n,r+1),\qquad B_P=(kH)^k.
\]
The positive alternative is equivalent to feasibility of
\begin{equation}
 C_Ix=0,\qquad c_I^{\mathsf T}x=1,
 \qquad 0\le x_j\le B_P.
 \label{eq:polynomial-lp-primal-box}
\end{equation}
Indeed, a positive kernel point rescales to the unbounded slice.  A nonempty
standard-form polyhedron has a basic feasible point, supported on at most
$\operatorname{rank}[C_I;c_I^{\mathsf T}]\le k$ columns.  Cramer's rule writes
its coordinates as ratios of integer determinants, with a nonzero integer
denominator of magnitude at least one.  The numerator is bounded by
$k!H^k\le B_P$, since the right side consists of zeroes and a one.  Thus the
box preserves at least one feasible point.

If this slice is empty, the homogeneous Farkas alternative supplies
$C_I^{\mathsf T}y_I\ge c_I$.  Because the selected rows are independent,
$C_I^{\mathsf T}$ is injective: this dual polyhedron has no lineality.
A nonempty pointed polyhedron has a vertex.  At such a vertex $r$ independent
inequalities are tight, and Cramer's rule bounds each coordinate by
$r!H^r\le B_D=(rH)^r$.  It therefore suffices to find a point in
\begin{equation}
 -C_I^{\mathsf T}y_I\le-c_I,
 \qquad -B_D\le (y_I)_i\le B_D.
 \label{eq:polynomial-lp-dual-box}
\end{equation}
For $r=0$ this is the zero-dimensional constant-inequality problem; the
implementation uses the radius convention $(\max(r,1)H)^r$.
The original-row answer is
\[
 y_{{\rm original},j_i}=\frac{s_i}{s_c}(y_I)_i,
 \qquad y_{{\rm original},j}=0\quad(j\notin\{j_1,\ldots,j_r\}).
\]
Consequently $C^{\mathsf T}y_{\rm original}\ge c$.  Importantly, a failed
primal decision does not by itself become a negative certificate: the dual
point is constructed separately and its original inequalities are checked.
For $n=0$ a zero multiplier is immediate.  For rank zero and $n>0$, a positive
objective coefficient gives a coordinate ray; otherwise the zero multiplier
works.  These cases also fit the exact zero/one-dimensional routines below.

\subsection{A determinant gap for exact bounded feasibility}

Consider an explicit bounded system
\[
 P=\{z\in\mathbb R^d:Az\le b\},
\]
including all box inequalities $-B\le z_j\le B$, where $B\ge0$ is rational.
Integerize each entire inequality, including its right side, by a positive
multiplier; common integer factors may then be removed.  Let
\[
 H_A=\max(1,|A_{ij}|),\qquad
 D=((d+1)H_A)^{d+1},\qquad \epsilon=\frac1{2D},
 \qquad P_\epsilon=\{z:Az\le b+\epsilon\mathbf1\}.
\]
The coefficient height $H_A$ need not include the right sides.  Their bit
length still contributes to the overall input size and to arithmetic costs.
A zero row with negative right side is an immediate contradiction.

\begin{proposition}[Exact perturbation gap]
The original system $P$ is feasible if and only if $P_\epsilon$ is feasible.
\end{proposition}
\begin{proof}
If $P$ is empty, Farkas multipliers can be normalized to satisfy
$\lambda\ge0$, $A^{\mathsf T}\lambda=0$, $\mathbf1^{\mathsf T}\lambda=1$,
and $b^{\mathsf T}\lambda<0$.  Minimize $b^{\mathsf T}\lambda$ on this
compact multiplier polytope.  A minimizing vertex uses at most $d+1$
independent columns of $[A^{\mathsf T};\mathbf1^{\mathsf T}]$.  Its
coordinates have a common nonzero integer determinant denominator with
magnitude at most $D$.  Since $b$ is integral, its negative objective is
at most $-1/D$.  Hence
\[
 \lambda^{\mathsf T}(b+\epsilon\mathbf1)
 \le-\frac1D+\frac1{2D}<0,
\]
which certifies that $P_\epsilon$ is empty too.  The reverse implication
uses $P\subseteq P_\epsilon$.
\end{proof}

If $z_0\in P$ and $d\ge1$, then $P_\epsilon$ contains the Euclidean ball
around $z_0$ of radius
\[
 r_0=\frac{\epsilon}{dH_A},
\]
because every row has norm at most $\sqrt d H_A\le dH_A$.  Each integerized
box row has coefficient magnitude at least one.  Thus
$P_\epsilon\subseteq[-B-1,B+1]^d$, which is enclosed by the radius
\[
 R=d(B+1)
\]
ball about zero.  This supplies a uniform positive inner radius even when
the original polytope is lower dimensional, as happens with matching
constraints and forced zero coordinates.  In dimension zero one checks
constant inequalities.  In dimension one exact interval intersection
solves the original system directly.  The ellipsoid construction below
is used only for $d\ge2$.

\subsection{Rational central cuts and explicit outward rounding}

Write
\[
 E(a,Q)=\{z:(z-a)^{\mathsf T}Q^{-1}(z-a)\le1\},
 \qquad Q\succ0.
\]
Start with $a=0$ and $Q=R^2I$.  At a rational center, scan the inequalities
of $P_\epsilon$ exactly.  If none is violated, the perturbation proposition
proves feasibility of the original system, although this center is not yet
an original feasible point.  Otherwise a nonzero violated row $g$ supplies
a central halfspace $g^{\mathsf T}(z-a)\le0$ containing $P_\epsilon$.
Let $w=Qg$ and $s=g^{\mathsf T}Qg>0$.  The ideal update is
\begin{align*}
 a_*&=a-\frac{w}{(d+1)\sqrt s},\\
 Q_*&=\frac{d^2}{d^2-1}
       \left(Q-\frac{2ww^{\mathsf T}}{(d+1)s}\right).
\end{align*}
It contains the retained half-ellipsoid.  To check this directly, map the
old ellipsoid to the unit ball and the cut to its first coordinate
$u\in[-1,0]$.  The new squared axes are $d^2/(d+1)^2$ in the cut direction
and $d^2/(d^2-1)$ perpendicular to it.  The new squared norm of a retained
unit-ball point is at most
\[
 1+\frac{2(d+1)u(u+1)}{d^2}\le1.
\]
For the ideal volume ratio $v_d$, putting $x=1/d$ gives
\[
 -2\log v_d
 =(1/x+1)\log(1+x)+(1/x-1)\log(1-x)
 =x+\sum_{j\ge1}\frac{x^{2j+1}}{(2j+1)(j+1)}\ge x.
\]
In particular $\log v_d\le-1/(2d)$.

Set the following explicit rational/integer parameters:
\begin{align*}
 K&=\lceil\log_2(R/r_0)\rceil,&
 N&=4d^2(K+1),& \rho&=\frac1{100d^2},\\
 p_i&=2(i+1)+\lceil\log_2(200d^3)\rceil,&
 \delta_i&=2^{-p_i},& \lambda_i&=4^{-(i+1)}.
\end{align*}
At update $i$, round every coordinate of $a_*$ down to the
$\delta_i$ grid to obtain $\widehat a$.  Form $Q_0=(1+\rho)^2Q_*$,
round each entry down to the same grid, and add $d\delta_i$ to each diagonal
entry.  Call the resulting symmetric matrix $\widehat Q$.  All logarithm
ceilings are computed with integer bit lengths and exact comparisons.

\begin{proposition}[Containment and progress of the stored update]
At every update the stored ellipsoid contains $P_\epsilon$.  Its shape is
positive definite, with minimum eigenvalue at least $R^2 4^{-(i+1)}$ after
update $i$.  Each stored update decreases log volume by more than $1/(3d)$.
\end{proposition}
\begin{proof}
The Cauchy--Schwarz inequality in the $Q$ inner product gives
$ww^{\mathsf T}/s\preceq Q$.  Consequently
\[
 Q_*\succeq\left(\frac d{d+1}\right)^2Q\succeq Q/4.
\]
Inductively $\lambda_{\min}(Q_*)\ge R^2\lambda_i\ge\lambda_i$, since
$R\ge1$.  The chosen grid satisfies
$\delta_i\le\rho\lambda_i/(2d)$.  Its center error has Euclidean norm at
most $\sqrt d\delta_i$, hence $Q_*$ norm at most
$\sqrt d\delta_i/\sqrt{\lambda_i}\le\rho$.  The norm triangle inequality
therefore gives
\[
 E(a_*,Q_*)\subseteq E(\widehat a,(1+\rho)^2Q_*).
\]
For the symmetric entrywise matrix-rounding error $T$, every entry has
absolute value at most $\delta_i$, so $\|T\|_2\le d\delta_i$.  It follows that
\[
 0\preceq T+d\delta_i I\preceq2d\delta_i I\preceq\rho Q_0,
 \qquad Q_0\preceq\widehat Q\preceq(1+\rho)^3Q_*.
\]
This proves positive definiteness, the claimed lower eigenvalue bound,
and outward containment.  The inflation changes volume relative to the
ideal update by at most $(1+\rho)^{3d/2}$.  Therefore
\[
 \log(\text{stored volume ratio})
 \le-\frac1{2d}+\frac{3d\rho}2
 =-\frac{97}{200d}<-\frac1{3d}.
\]
\end{proof}

After $N$ updates without an accepted center, the enclosing volume is less
than $e^{-N/(3d)}$ times the initial volume.  The definition of $N$ gives
$N/(3d)>d\log(R/r_0)$, so this is smaller than the volume of the inner-radius
ball forced by any feasible original point.  Reaching this cutoff therefore
proves infeasibility.  It is an exact volume argument, not an iteration
heuristic.  Plain rational arithmetic without the explicit grid, inflation,
and diagonal correction would not prove these claims.

\subsection{No irrational arithmetic and no unbounded stored precision}

The adaptive grids are nested, with precision increasing by two each update.
The old center is therefore integral after multiplication by $2^{p_i}$.
The required rounded center coordinate is exactly
\[
 2^{p_i}\widehat a_j
 =2^{p_i}a_j-\left\lceil\frac{t_j}{\sqrt s}\right\rceil,
 \qquad t_j=\frac{2^{p_i}w_j}{d+1}.
\]
For $t_j>0$, set $u=t_j^2/s$.  Then
$\lfloor t_j/\sqrt s\rfloor=\operatorname{isqrt}
(\lfloor\operatorname{num}(u)/\operatorname{den}(u)\rfloor)$.
The exact square comparison decides whether the ceiling needs one more.
For negative $t_j$, the required ceiling is the negative of the corresponding
nonnegative floor; zero is immediate.  Thus rational operations and integer
square root implement the ideal-center rounding exactly.

The upper shape inequality is also explicit:
\[
 \widehat Q\preceq(1+\rho)^3\frac{d^2}{d^2-1}Q\preceq2Q
 \quad(d\ge2).
\]
The largest shape eigenvalue is at most $R^2 2^i$.  An ideal center move has
Euclidean norm at most $\sqrt{\lambda_{\max}(Q)}/(d+1)$, followed by a grid
error of at most $\sqrt d\delta_i$.  Hence center magnitudes have
$O(N+\log R+\log N+\log d)$ bits before denominator accounting.
After an update every shape and center entry is dyadic with denominator
$2^{p_i}$, and
\[
 p_i\le2N+\lceil\log_2(200d^3)\rceil=O(N+\log d).
\]
The initial rational shape need not be dyadic; its encoding is bounded by
that of $R^2$.  Exact intermediate formulas combine only polynomially many
bits: the products $Qg$ share the dyadic shape denominator, the row $g$ is
integer, and the scalar $s$ supplies an explicit rational denominator for
the rank-one update.  The integer-square-root inputs obey the same bound.
Since $\log D=O(d\log(dH_A))$,
\[
 N=O\bigl(d^2(\log R+d\log(dH_A))\bigr).
\]
Both the number of cuts and their bit cost are polynomial.  Starting with
the final precision on every update is unnecessary; the finer successive
grids retain this worst-case proof without imposing late precision early.

\subsection{Recovering an exact original point by tightening rows}

An accepted ellipsoid center solves only the perturbed inequalities.  To
recover an original feasible point, first decide that $P$ is nonempty by
the exact decision procedure above.  Maintain a nonempty face $F$, initially
$P$.  Scan the original inequalities, including the box rows.  For
$a_i^{\mathsf T}z\le b_i$, decide feasibility of
\[
 F\cap\{z:a_i^{\mathsf T}z\ge b_i\}.
\]
When feasible, retain that reverse inequality permanently, making the
original row tight.  Each such decision uses the same determinant-gap
procedure, applied to the current explicit row system.

After the scan the final face is a singleton.  Otherwise a positive-dimensional
bounded polytope $F$ has a proper nonempty boundary face where some original
inequality becomes tight.  When that row was inspected, the then-current
face contained this final boundary face, so the equality test was feasible.
The reverse inequality would have been retained, contradicting that the row
is nonconstant on final $F$.  Thus the retained equalities have rank $d$.
Exact Gaussian elimination obtains their unique point, and every original
inequality is checked directly.  One may stop as soon as rank $d$ is reached.

There are at most one plus the original number of rows decision calls.
Reverse inequalities have exactly the original integer coefficient and
right-side heights, and the row count at most doubles.  Thus this recovery
is polynomial too.  It does not guess an exact answer by rounding a nearly
feasible center, and it requires no simultaneous Diophantine reconstruction.
An exhausted shared decision/cut allowance returns no partial point.

\subsection{From a boxed point to an extreme kernel ray}

The recovered point of~\eqref{eq:polynomial-lp-primal-box} has positive
objective, but the artificial box can create vertices that are not kernel
extreme rays.  This distinction matters: the uniform resolving-set theorem
on positive inputs uses root extreme-ray conflicts.  The LP wrapper therefore
runs the generic original-source support descent on the recovered point,
using $C$ and $c$ with an \emph{empty} compatibility list.

Normalize $\mathbf1^{\mathsf T}x=1$.  If the matching kernel on its positive
support has dimension greater than one, its normalized face admits a nonzero
tangent direction $d$ with $Cd=0$ and $\mathbf1^{\mathsf T}d=0$.  Orient $d$
so $c^{\mathsf T}d\ge0$, move to the first coordinate-zero endpoint, and
retain all previous zeroes.  Such an endpoint exists in the bounded simplex;
it deletes a positive coordinate while preserving nonnegative matching and
strictly positive objective.  At most $n-1$ steps reach a support of matching
nullity one.  The final primitive integer vector is an extreme ray of the
original nonnegative kernel cone.

This procedure also has an explicit bound on its \emph{intermediate} bit
lengths.  Clear the original matching denominators and let $H\ge1$ bound
the resulting integer height.  The source identity chart describes each
chosen direction by original matching rows, accumulated coordinate-zero
rows, the row of ones, and unit constraints on surviving original free
coordinates.  The chart's parameters are actual original source coordinates;
its restrictions preserve the identity submatrix at the surviving free
indices.  Selecting one chart direction therefore amounts to setting one
such original free coordinate to one and the others to zero in this
original-variable system.  Select a nonsingular square subsystem.  Its
integer coefficient height is still at most $H$, including the unit and
sum rows, and its right side has only zeroes and ones.  All rational
direction coordinates have a common Cramer denominator; clearing that
single denominator and removing a gcd bounds every primitive integer
entry by
\[
 M=n!H^n
\]
Write the current normalized point as $v/D$ with integral
$v$ and $D>0$.  If $d_j<0$ is the coordinate killed at the endpoint, then
\[
 x'_k=\frac{v_k(-d_j)+v_jd_k}{D(-d_j)}.
\]
A common denominator is multiplied by at most $M$ per step, so
$\log D_{\rm final}\le\log D_{\rm initial}+n\log M$.
All intermediate coordinates remain in the unit simplex, bounding their
numerators as well.  The initial recovered vertex has determinant-bounded
rational coordinates; taking a common denominator and normalizing its sum
adds only polynomially many bits.  This justifies the actual support-descent
trajectory, not merely the final ray's encoding length.

For the legacy basis output, let $C_T$ denote the rows retained in the
simplex tableau, and consider its slack matrix
\[
 \left[\begin{array}{c} C_T\\-C_T\\\mathbf1^{\mathsf T}\end{array}
 \ \middle|\ I\right].
\]
The positive support columns are independent after the normalization row
is appended, since their matching nullity is one.  Extend these columns
to a basis using the remaining columns, which is always possible because
of the identity slack block.  The normalized ray with all matching and
normalization slacks zero is the corresponding nonnegative basic solution.
The fallback independently checks the support rank and basis rank, as well
as the original-source primal inequalities.  When the explicit
\texttt{reduce\_rows=False} ablation is used, $C_T$ includes every original
row for legacy slack indexing, while the bounded feasibility systems still
use an independent original-row subset.

This generic kernel extraction is separate from the later extraction in
\emph{original standard normal coordinates} available for arbitrary positive
points before crushing.  For the native primitive Q extreme outputs,
canonical lifting already gives a primitive STANDARD extreme ray;
Section~\ref{sec:direct-crush-positive} removes that redundant later pass
while retaining independent original-support and quotient checks.  An
unrestricted Q ray may be incompatible.  Neither its positive linear
objective nor its extremality alone asserts an embedded surface: the Euler
oracle still checks the actual canonical formal Euler and the quadrilateral
condition before any positive normal-surface result is returned.

\subsection{The Bland prefix has polynomial bit cost}

Let $r_T$ be the number of tableau rows selected from the input, and put
\[
 P_0=(n+r_T+1)^2.
\]
The polynomial policy allows at most $P_0$ Bland pivots as its fast path.
The optional \path{fast_path_pivots} can only lower this bound.
A conclusive original simplex answer is returned immediately;
otherwise the internal cap triggers the exact fallback.  The explicit
\texttt{policy='bland'} ablation retains the old finite but potentially
exponential uncapped simplex trajectory.

Bounding the number of pivots is only part of the argument.  Put
$m_T=2r_T+1$ and integerize the \emph{fixed} source slack system.  Row clearing
scales its slack coefficient and right side along with its source
coefficients.  Accordingly the integer height $H_T$ must include the
row-clearing multipliers, not just the integerized matching entries.
At every pivot the basis matrix is a nonsingular $m_T$-column submatrix of
this fixed system.  The tableau and right side are $B^{-1}A$ and $B^{-1}b$.
Cramer's rule bounds their reduced numerators and denominators by minors of
the fixed input matrix, giving
\[
 O\bigl(m_T(\log m_T+\log H_T)\bigr)
\]
bits per entry.  Clearing objective denominators gives a polynomial bound
for $c_j-c_B^{\mathsf T}B^{-1}A_j$ and $c_B^{\mathsf T}B^{-1}b$ as well.
Every pivot performs polynomially many rational additions, multiplications,
and divisions.  Exact Fraction reduction after each operation bounds the
temporary products and sums polynomially.  Thus the $P_0$-pivot prefix has
polynomial bit cost even along a degenerate simplex sequence.  Invoking
Bland finiteness alone would not prove this.

Independent-original-row selection likewise uses an exact echelon chart
with entries given by ratios of source minors.  The controller shares strict
parsing and this selection with the fast path.  Its conclusive results keep
the established basic-basis invariant and original primal/dual checks;
additional independent support/basis elimination is reserved for fallback
answers and research replay.  This avoids repeating fallback-only work on
every ordinary fast-path result without weakening its mathematical contract.
Section~\ref{sec:prepared-euler-matrix} additionally shares this immutable
matrix preparation across the uncached objectives of one Euler mask,
while retaining fresh objective-dependent solver state and allowances.

\subsection{A conservative explicit binary-complexity bound}

\begin{proposition}[Polynomial kernel-alternative policy]
On explicit rational input, with no finite caller caps and constant-cost
cancellation polling, the specified polynomial policy returns a positive
original-cone extreme ray or an original-row dual in polynomial bit time.
\end{proposition}
\begin{proof}
The capped simplex prefix has polynomial bit cost by the preceding argument.
The primal and dual boxes preserve feasible witnesses of the relevant
alternative.  The determinant gap, outward rounding, and volume cutoff give
an exact polynomial bounded-feasibility decision.  Tight-row recovery uses
polynomially many such decisions.  Source support descent and slack-basis
completion are polynomial and preserve the required extreme-ray contract.
The separate final original-system checks establish the claimed alternatives.

For an explicit deliberately loose bound, let $L$ be the total binary length
of all rational input coefficients, and set
\[
 U=2+n+m+L.
\]
Denominator clearing gives $\log H=O(U)$ and primal/dual box radii with
$O(U^2)$ bits.  Each fallback system has dimension and row count $O(U)$.
One decision needs $O(U^4)$ cuts, each exact recovery needs $O(U)$ decisions,
and the whole primal/dual fallback therefore needs $O(U^5)$ cuts.  Each cut
uses $O(U^2)$ rational/rounding operations.  The adaptive stored entries and
intermediate quantities have $O(U^4)$ bits; the retained-face solves,
source descent, and polynomial Bland prefix fit within these loose bounds.
Thus $O(U^7)$ exact arithmetic or integer-square-root operations on
$O(U^4)$-bit operands suffice.  Even elementary cubic bit bounds for those
operations give the conservative bound $O(U^{19})$.  Its purpose is to make
the polynomial claim explicit, not to describe expected practical speed.
\end{proof}

The bound includes rational input magnitudes and denominator lengths; it is
not a strongly polynomial bound.  Faster integer arithmetic and tighter rank
accounting improve it.  Work performed inside an arbitrary caller callback
is outside the solver's own complexity bound.

\subsection{Finite allowances, cancellation, and the remaining global gap}

An explicit \texttt{max\_pivots} remains a hard caller allowance.  If it is
exhausted before a conclusive answer, the wrapper returns
\texttt{INCONCLUSIVE}; fallback cannot bypass it.  Reaching only the internal
polynomial prefix limit starts fallback.  Separate
\texttt{max\_ellipsoid\_iterations} and \texttt{max\_feasibility\_calls}
allowances count every fallback update and every feasibility decision,
including all primal recovery and dual recovery calls.  No ellipsoid update
is described as a simplex pivot.  Exact constant/interval decisions may
finish with zero updates, but still count as decisions.  Exhausted caps do
not produce a primal, dual, or basis payload.

Cancellation polling occurs in preparation, row scans, rounding/shape loops,
equality recovery, and support descent.  The callback is invoked directly:
its exception object propagates unchanged, including exceptions normally used
for validation and subclasses of \texttt{BaseException}.  There is no broad
exception-to-infeasibility conversion.  One integer or rational operation
is not independently preemptible.

The maintained Euler callers share these counters across objectives,
anchors, masks, candidate restarts, and authenticated crushing successors.
The common \path{euler_lp_budget.py} helper subtracts previously charged
work before each LP and merges all returned counters, including capped
attempts.  The compatible-sector, tree, and transversal APIs expose
\path{lp_policy='polynomial'}, \path{lp_fast_path_pivots},
\path{max_ellipsoid_iterations}, and \path{max_feasibility_calls}.
Their direct defaults leave fallback work uncapped; the public diagram
crushing driver defaults the two fallback allowances to $10000$ each.
Passing \texttt{None} disables an allowance.  The explicit Bland comparison
ignores unused fallback budgets at these caller interfaces, while the
low-level solver rejects fallback-only controls under its Bland policy.

The flat fair sector scheduler applies its local attempt quota separately
to pivots, anchors, ellipsoid updates, and feasibility calls, in addition
to their global remaining allowances.  Tree and transversal search share
one oracle budget throughout their nodes and candidate restarts; the
crushing driver subtracts every prior successor's work before starting
another source.  A fixed-source dual-cache hit needs no new LP call;
uncached calls include their entire capped prefix and fallback.  Summed
statistics distinguish \path{lp_fallbacks}, \path{lp_ellipsoid_iterations},
\path{lp_feasibility_calls}, and \path{lp_support_descent_steps}; maximum
precision, dimension, and planned-cut statistics are merged by maximum.
Independent negative-tree replay continues to check original source dual
inequalities and does not rerun either simplex or the ellipsoid algorithm.

The integrated regression run records 192 tests in $55.862$ seconds,
with one optional Regina census test skipped because the system interpreter
lacks Regina.  That single test was rerun under the Regina interpreter and
passed in $0.178$ seconds.  The logs are
\path{data/exact-polynomial-lp-tests.txt} and
\path{data/exact-polynomial-lp-regina-test.txt}.  New caller-budget controls
exercise genuine one-dimensional recovery, real fallback work across
transversal restarts, a shared cut cap across fair retries, and complete
fast paths with unused zero fallback allowances.  The successor-accounting
control uses explicitly synthetic discovery costs on three genuine crushing
steps; it tests budget propagation rather than measuring discovery.  These
regressions support the implementation contract.  They are neither an
isolated performance experiment nor a substitute for the bit-complexity
proof, and no fallback speedup is claimed here.

For a prepared source with $A$ canonical Euler anchors, this policy supplies
polynomial bit work per exact LP call.  Conditional on a resolving number
$h$, the transversal argument therefore retains its bound
\[
 A(3t)^{O(h)}\operatorname{poly}(t,L),
\]
with the former unbounded-Bland hypothesis discharged for this policy.
After one-vertex preprocessing the source anchor count is at most $4t$.
The unresolved structural obligation is a uniform small resolving parameter,
or a different bound on the complete admissibility search.  General branch
coverage remains exponential; the finite $h=3$ and $h=2$ examples do not bound
a family.  Consequently polynomial exact LP, even together with certified
polynomially many crushing stages, does not establish quasi-polynomial
unknot recognition.

\subsection{Fresh exact audit of the maintained policy}

The source-pinned package \path{fast/polynomial_lp_research} performs fresh
discovery with the maintained implementation.  Its 15 rational kernel cones
are each queried under both row policies and both the explicit Bland policy
and the polynomial policy with a zero-pivot prefix.  The independent oracle
enumerates the normalized cone vertices using Leibniz determinants and
Cramer's rule, without any producer elimination helper.  All 60 queries
agree: 20 are positive and 40 nonpositive.  Every positive answer is checked
for original matching, primitive normalization, support nullity one, and a
nonsingular original slack basis.  Every negative answer satisfies all
original-row dual inequalities.  Twenty-two queries actually enter fallback,
using 16642 ellipsoid updates and 100 feasibility decisions in total;
trivial zero-pivot conclusions are not counted as fallback discoveries.

The inputs include dependent rational rows, unsorted pivots, a genuine
four-dimensional negative problem, and seeded small controls.  Two further
three-dimensional bounded systems have independently known outcomes: an
explicit singleton and a pair of inequalities whose sum is a strict
constant contradiction.  Twelve consecutive updates from each are retained.
Their solver-free replay checks all principal minors for outward matrix
containment, positive definiteness, the spectral floor, a rational volume
decrease bound, nested dyadic precision, and exact center rounding.  It also
checks 51 signed integer-square-root ceilings by squared inequalities.
These 24 recorded updates test implementation arithmetic; the general
containment and cutoff claims come from the proof above.

Two genuine geometric queries force fallback as well.  A two-tetrahedron
layered-torus mask gives a complete negative tree after candidate restarts;
its source-derived dual leaves and complete coverage replay independently.
A one-tetrahedron source is obtained from an authenticated partial crushing
chain of the circle.  Its positive sector is freshly solved by fallback,
checked in original normal coordinates, extracted, and used to complete an
independently replayed crushing chain.  The small-source tests establish
neither an exhaustive family audit nor a practical advantage for forcing
the ellipsoid path.

The saved evidence is \path{data/polynomial-lp-audit.json}.  The independent
checker runs with \texttt{python -S}, invokes no LP or crushing discovery,
and replays all retained original-system and geometric certificates.  The
audit and timing files pin every native Python module and their precise
research dependencies; the final publication manifest additionally pins
the article, compiled PDF, regression logs, and generated timing table.

\subsection{Complete-query comparison with the retained Bland policy}

The separate paired experiment compares \path{lp_policy='polynomial'} with
\path{lp_policy='bland'} in the same pinned native implementation.  Both
use source-model transversal search, envelope-first objectives, vertex
reduction, and identical finite shared allowances: 256 tree nodes and
10000 each for pivots, anchors, tested subsets, ellipsoid updates, and
feasibility decisions.  Each case has three shuffled paired samples and
one shuffled Bland/Bland A/A pair, with batch size one.  The fixed seed is
2026101301.  Construction, shelling, preprocessing, discovery, internal
replay, and full result serialization are timed; diagram objects, imports,
and a separate independent warmup replay are outside the timer.

\begin{center}
\begin{tabular}{lrrrr}
\toprule
Case & Bland (ms) & Polynomial (ms) & Bland/polynomial & A/A \\
\midrule
circle & 9.655 & 9.112 & 0.997 & 1.033 \\
trefoil & 3088.487 & 3315.942 & 0.985 & 1.017 \\
figure\_eight & 9350.999 & 9384.090 & 0.998 & 0.999 \\
\bottomrule
\end{tabular}

\end{center}
The ratio column is the median of the paired ratios, not the ratio of
the two marginal medians.  Every timed output equals its independently
replayed warmup byte for byte, and the full outputs also agree between
policies.  Both policies complete the circle as \texttt{UNKNOT} and the
trefoil and figure-eight as \texttt{KNOTTED}.  No query in this experiment
enters ellipsoid fallback.  Vertex preprocessing alone resolves the circle,
so that case performs no LP call; the trefoil and figure-eight use 124 and
114 LP calls, with 1811 and 5433 pivots, respectively.

The paired Bland/polynomial ratios are $0.997$, $0.985$, and $0.998$;
the corresponding A/A ratios are $1.033$, $1.017$, and $0.999$.
These short measurements support retaining the common fast path and expose
its finite overhead and timing variation.  They establish neither a
practical fallback speedup nor a broad performance claim.  The improvement
proved here is the worst-case binary complexity of each uncapped LP query.
Raw samples, complete warmups, source pins, and independently recomputed
table values are in \path{data/polynomial-lp-benchmark.json} and
\path{data/polynomial-lp-replay.txt}.  Reproduction commands are in
\path{fast/polynomial_lp_research/README.md}.
